
\begin{table}[!htbp]
\centering
\caption{Perbandingan pemenuhan kebutuhan sistem pada empat kondisi penelitian.}
\label{tab:perbandingan-kondisi}
\small
\renewcommand{\arraystretch}{1.2}

\begin{tabularx}{\textwidth}{|>{\raggedright\arraybackslash}X|c|c|c|c|}
\hline
\textbf{Kriteria Kebutuhan}
  & \textbf{\makecell{Kondisi A\\LLM Murni}}
  & \textbf{\makecell{Kondisi B\\RAG\\Konvensional}}
  & \textbf{\makecell{Kondisi C\\GraphRAG\\Berbobot}}
  & \textbf{\makecell{Kondisi D\\GraphRAG\\\emph{Dual-Level}}}
\\ \hline

\emph{Grounding} arsitektural berbasis Stack Overflow (G1)
  & $\times$ & $\sim$ & $\checkmark$ & $\checkmark$ \\ \hline
Pengetahuan \emph{library}/API terkini (G1)
  & $\times$ & $\checkmark$ & $\checkmark$ & $\checkmark$ \\ \hline
Sinyal kepercayaan komunitas Stack Overflow (G2)
  & $\times$ & $\times$ & $\checkmark$ & $\times$ \\ \hline
\emph{Retrieval} relasional dan multi-hop (G2, G3)
  & $\times$ & $\times$ & $\checkmark$ & $\checkmark$ \\ \hline
\emph{Semantic retrieval} berbasis vektor
  & $\times$ & $\checkmark$ & $\checkmark$ & $\checkmark$ \\ \hline
\emph{Grounded generation} dan sitasi (G1)
  & $\times$ & $\sim$ & $\checkmark$ & $\checkmark$ \\ \hline
\emph{Provenance} dan keterlacakan sumber
  & $\times$ & $\sim$ & $\checkmark$ & $\sim$ \\ \hline
Protokol evaluasi komparatif
  & $\checkmark$ & $\checkmark$ & $\checkmark$ & $\checkmark$ \\ \hline
Pembaruan pengetahuan tanpa \emph{re-training}\textsuperscript{\dag}
  & $\times$ & $\sim$ & $\checkmark$ & $\checkmark$ \\ \hline

\end{tabularx}

\vspace{0.2cm}
\footnotesize
$\checkmark$ = Terpenuhi \qquad
$\sim$ = Sebagian terpenuhi \qquad
$\times$ = Tidak terpenuhi \\
Kondisi B memperoleh $\sim$ pada \emph{grounding} dan sitasi karena instruksi tersebut bersifat opsional. Kondisi D memperoleh $\sim$ pada \emph{provenance} karena konteksnya tidak membawa bobot kepercayaan. \\
\textsuperscript{\dag}\todo{Baris ini bergantung pada keputusan F7.}

\end{table}
